\documentclass[12pt]{article}

\usepackage{amsmath}
\usepackage{graphicx}
\usepackage{booktabs}
\usepackage[a4paper, margin=1in]{geometry}

\title{Labwork 1: Gradient Descent}
\author{Nguyen Xuan Vinh}
\date{\today}

\begin{document}

\maketitle

\section{Introduction}

Gradient Descent is one of the most fundamental optimization algorithms in Machine Learning and Deep Learning. 
It is used to iteratively minimize a function by updating parameters in the direction opposite to the gradient.

In this labwork, Gradient Descent was implemented from scratch using Python to minimize the following function:

\[
f(x) = x^2
\]

The derivative of the function is:

\[
f'(x) = 2x
\]

The update rule of Gradient Descent is:

\[
x_{new} = x - r f'(x)
\]

where:
\begin{itemize}
    \item \(x\) is the current value
    \item \(r\) is the learning rate
    \item \(f'(x)\) is the derivative of the function
\end{itemize}

The objective is to find the minimum value of the function.

\section{Algorithm Implementation}

The algorithm was implemented manually without using any optimization libraries.

The implementation process is summarized as follows:

\begin{enumerate}
    \item Initialize a starting value \(x\)
    \item Compute the gradient \(f'(x)\)
    \item Update \(x\) using the Gradient Descent formula
    \item Repeat the process for a fixed number of iterations
    \item Print intermediate results at each iteration
\end{enumerate}

The Python program contains:
\begin{itemize}
    \item A function to compute \(f(x)\)
    \item A function to compute the derivative \(f'(x)\)
    \item A Gradient Descent loop for iterative optimization
\end{itemize}

\section{Experimental Results}

The initial value was set to:

\[
x = 10
\]

Different learning rates were tested to analyze their effects on convergence.

\subsection{Learning Rate \(r = 0.1\)}

This learning rate produces stable convergence toward the minimum point. 
However, the convergence speed is relatively slow because each update step is small.

\begin{table}[h]
\centering
\begin{tabular}{ccc}
\toprule
Iteration & x & f(x) \\
\midrule
0 & 10.0000 & 100.0000 \\
1 & 8.0000 & 64.0000 \\
2 & 6.4000 & 40.9600 \\
3 & 5.1200 & 26.2144 \\
4 & 4.0960 & 16.7772 \\
\bottomrule
\end{tabular}
\caption{Gradient Descent with learning rate \(r = 0.1\)}
\end{table}

\subsection{Learning Rate \(r = 0.5\)}

This learning rate converges significantly faster and reaches the minimum efficiently. 
The updates are larger but still stable.

\subsection{Learning Rate \(r = 1.1\)}

This learning rate is too large and causes unstable behavior. 
The algorithm overshoots the minimum point repeatedly, leading to divergence and oscillation.

\section{Discussion}

The experiments demonstrate that the learning rate is a critical hyperparameter in Gradient Descent.

\begin{itemize}
    \item A small learning rate leads to slow but stable convergence.
    \item A suitable learning rate improves optimization efficiency.
    \item An excessively large learning rate causes divergence and instability.
\end{itemize}

Therefore, selecting an appropriate learning rate is important for successful optimization in Machine Learning models.

\section{Conclusion}

In this labwork, Gradient Descent was successfully implemented from scratch to minimize the function:

\[
f(x) = x^2
\]

The implementation correctly updates the variable iteratively using the derivative of the function.

The experimental results show that different learning rates strongly affect convergence speed and stability. 
This labwork provides a practical understanding of how Gradient Descent works in optimization problems.

\end{document}